% انتخاب کلاس سند
\documentclass[12pt,a4paper]{article}

% ---------- پکیج‌ها (همه قبل از xepersian) ----------
\usepackage{fontspec}
\usepackage{amsmath, amssymb, amstext}
\usepackage{graphicx}
\usepackage{geometry}
\usepackage{xcolor}
\usepackage{fancyhdr}

% ---------- تنظیم حاشیه ----------
\geometry{margin=1in}

% ---------- تنظیم fancyhdr ----------
\pagestyle{fancy}
\setlength{\headheight}{15pt} % حل خطای headheight

\lhead{کوییز درس مکانیک تحلیلی}
\chead{دانشکده فیزیک دانشگاه صنعتی خواجه نصیرالدین طوسی}
\rhead{\today}

% ---------- فارسی (باید آخر بیاید) ----------
\usepackage{xepersian}
\settextfont{B Nazanin}

% ---------- شروع سند ----------
\begin{document}
	
	\begin{tabular}{|p{7cm}|p{7cm}|}
		\hline
		\textbf{نام و نام خانوادگی:} & \textbf{شماره دانشجویی:} \\
		\hline
	\end{tabular}
	
	\section*{سوالات:}
	
	\begin{enumerate}
		\item 
		سطحی دوار را در نظر بگیرید که از دوران منحنی \( y(x) \) حول محور \( x \)
		ایجاد می‌شود.  
		این منحنی از دو نقطه انتهایی ثابت
		\[
		(x_1, y_1), \qquad (x_2, y_2)
		\]
		می‌گذرد.
		مسئله وردشی را به‌گونه‌ای فرموله کنید که تابع \( y(x) \) به‌عنوان تابع
		اختیار انتخاب شود و مساحت سطح دوار کمینه گردد.
		\item
		معادلات ماشین آتوود مرکب را بدست آورید 
		
	\end{enumerate}
	
	\vspace{10pt}
	\textbf{موفق باشید.}
	
\end{document}
